% GENERATED FILE. Edit canonical lessons, then run npm run generate:books.
% canonical-source-hash: fnv1a64:1c03429eaf05315a
% canonical-lessons: HI-C334-katha, HI-C334-rsi, HI-C334-asram, HI-C334-paheli, HI-C334-kahavat

\chapter{A Tale of Long Ago}
\label{ch:hi-a-tale-of-long-ago}
\hlchaptermodality{334}

\begin{chapteropening}
\textbf{By the end of this chapter:} \emph{I can begin and tell a traditional tale.}

Complete the last lesson of chapter 334: I can begin and tell a traditional tale.
\end{chapteropening}

\section[\texorpdfstring{\hi{कथा} (kathā)}{कथा (kathā)}]{\hi{कथा} (kathā) — a tale; a sacred story}

\label{lesson:HI-C334-katha}



\begin{quote}
Before the new one: say the Hindi for to be scattered, then the Hindi for the floating of lamps.
\end{quote}

\subsection*{You'll want to know: \hi{कथा}}
\textbf{\hi{कथा}} — \emph{kathā} — \textquotedblleft{}a tale; a sacred story\textquotedblright{}.

\begin{quote}
\hi{दादी} \hi{हर} \hi{रात} \hi{एक} \hi{नई} \hi{कथा} \hi{सुनाती} \hi{थीं।}
\end{quote}

\begin{quote}
\emph{dādī har rāt ek naī kathā sunātī thīṁ.}
\end{quote}

\begin{quote}
Grandmother used to tell a new tale every night.
\end{quote}

Feminine.

\begin{cousinweb}
\textbf{\hi{कथा}} (\emph{kathā}) is Sanskrit \emph{kathā}, \textquotedblleft{}a story, a conversation\textquotedblright{} (Monier-Williams). Beside \textbf{\hi{कहानी}}, it is the old or sacred tale.
\end{cousinweb}

\begin{grammarlens}[title={What you've built}]
One more word for telling what happened.
\end{grammarlens}

\subsection*{Your turn}
\begin{itemize}
\raggedright
  \item \emph{Say it:} \emph{kathā}
  \item \emph{Say it:} \emph{kathā}, once more
  \item \emph{Say it:} \emph{dādī har rāt ek naī kathā sunātī thīṁ.} — Grandmother used to tell a new tale every night.
  \item \emph{Recall:} say the Hindi for to be scattered, then the Hindi for the floating of lamps, then say \emph{kathā} again
\end{itemize}

\begin{tcolorbox}[breakable,colback=teal!4,colframe=teal!35!black,title={Before you move on}]
What does \hi{कथा} mean? (\textquotedblleft{}A tale; a sacred story\textquotedblright{}.) Say it once more.
\end{tcolorbox}
\section[\texorpdfstring{\hi{ऋषि} (ṛṣi)}{ऋषि (ṛṣi)}]{\hi{ऋषि} (ṛṣi) — a sage; a seer}

\label{lesson:HI-C334-rsi}



\begin{quote}
Before the new one: say the Hindi for the floating of lamps, then the Hindi for a tale.
\end{quote}

\subsection*{You'll want to know: \hi{ऋषि}}
\textbf{\hi{ऋषि}} — \emph{ṛṣi} — \textquotedblleft{}a sage; a seer\textquotedblright{}.

\begin{quote}
\hi{जंगल} \hi{में} \hi{एक} \hi{बूढ़े} \hi{ऋषि} \hi{रहते} \hi{थे।}
\end{quote}

\begin{quote}
\emph{jaṅgal meṁ ek būṛhe ṛṣi rahte the.}
\end{quote}

\begin{quote}
An old sage lived in the forest.
\end{quote}

Masculine.

\begin{cousinweb}
\textbf{\hi{ऋषि}} (\emph{ṛṣi}) is Sanskrit \emph{ṛṣi}, \textquotedblleft{}a seer, a sage\textquotedblright{} (Monier-Williams). The verb is plural for respect: \textbf{\hi{ऋषि} \hi{रहते} \hi{थे}}.
\end{cousinweb}

\begin{grammarlens}[title={What you've built}]
One more word for telling what happened.
\end{grammarlens}

\subsection*{Your turn}
\begin{itemize}
\raggedright
  \item \emph{Say it:} \emph{ṛṣi}
  \item \emph{Say it:} \emph{ṛṣi}, once more
  \item \emph{Say it:} \emph{jaṅgal meṁ ek būṛhe ṛṣi rahte the.} — An old sage lived in the forest.
  \item \emph{Recall:} say the Hindi for the floating of lamps, then the Hindi for a tale, then say \emph{ṛṣi} again
\end{itemize}

\begin{tcolorbox}[breakable,colback=teal!4,colframe=teal!35!black,title={Before you move on}]
What does \hi{ऋषि} mean? (\textquotedblleft{}A sage; a seer\textquotedblright{}.) Say it once more.
\end{tcolorbox}
\section[\texorpdfstring{\hi{आश्रम} (āśram)}{आश्रम (āśram)}]{\hi{आश्रम} (āśram) — a hermitage; an ashram}

\label{lesson:HI-C334-asram}



\begin{quote}
Before the new one: say the Hindi for a tale, then the Hindi for a sage.
\end{quote}

\subsection*{You'll want to know: \hi{आश्रम}}
\textbf{\hi{आश्रम}} — \emph{āśram} — \textquotedblleft{}a hermitage; an ashram\textquotedblright{}.

\begin{quote}
\hi{ऋषि} \hi{का} \hi{आश्रम} \hi{नदी} \hi{के} \hi{किनारे} \hi{था।}
\end{quote}

\begin{quote}
\emph{ṛṣi kā āśram nadī ke kināre thā.}
\end{quote}

\begin{quote}
The sage's hermitage was on the riverbank.
\end{quote}

Masculine.

\begin{cousinweb}
\textbf{\hi{आश्रम}} (\emph{āśram}) is Sanskrit \emph{āśrama}, \textquotedblleft{}a hermitage\textquotedblright{} (Monier-Williams): the forest home of a sage and his pupils.
\end{cousinweb}

\begin{grammarlens}[title={What you've built}]
One more word for telling what happened.
\end{grammarlens}

\subsection*{Your turn}
\begin{itemize}
\raggedright
  \item \emph{Say it:} \emph{āśram}
  \item \emph{Say it:} \emph{āśram}, once more
  \item \emph{Say it:} \emph{ṛṣi kā āśram nadī ke kināre thā.} — The sage's hermitage was on the riverbank.
  \item \emph{Recall:} say the Hindi for a tale, then the Hindi for a sage, then say \emph{āśram} again
\end{itemize}

\begin{tcolorbox}[breakable,colback=teal!4,colframe=teal!35!black,title={Before you move on}]
What does \hi{आश्रम} mean? (\textquotedblleft{}A hermitage; an ashram\textquotedblright{}.) Say it once more.
\end{tcolorbox}
\section[\texorpdfstring{\hi{पहेली} (pahelī)}{पहेली (pahelī)}]{\hi{पहेली} (pahelī) — a riddle}

\label{lesson:HI-C334-paheli}



\begin{quote}
Before the new one: say the Hindi for a sage, then the Hindi for a hermitage.
\end{quote}

\subsection*{You'll want to know: \hi{पहेली}}
\textbf{\hi{पहेली}} — \emph{pahelī} — \textquotedblleft{}a riddle\textquotedblright{}.

\begin{quote}
\hi{ऋषि} \hi{ने} \hi{राजा} \hi{से} \hi{एक} \hi{पहेली} \hi{पूछी।}
\end{quote}

\begin{quote}
\emph{ṛṣi ne rājā se ek pahelī pūchī.}
\end{quote}

\begin{quote}
The sage asked the king a riddle.
\end{quote}

Feminine.

\begin{cousinweb}
\textbf{\hi{पहेली}} (\emph{pahelī}) goes back to Sanskrit \emph{prahelikā}, \textquotedblleft{}a riddle\textquotedblright{} (CDIAL 8911). You \textbf{\hi{पहेली} \hi{पूछना}}, ask one, and \textbf{\hi{पहेली} \hi{बूझना}}, solve it.
\end{cousinweb}

\begin{grammarlens}[title={What you've built}]
One more word for telling what happened.
\end{grammarlens}

\subsection*{Your turn}
\begin{itemize}
\raggedright
  \item \emph{Say it:} \emph{pahelī}
  \item \emph{Say it:} \emph{pahelī}, once more
  \item \emph{Say it:} \emph{ṛṣi ne rājā se ek pahelī pūchī.} — The sage asked the king a riddle.
  \item \emph{Recall:} say the Hindi for a sage, then the Hindi for a hermitage, then say \emph{pahelī} again
\end{itemize}

\begin{tcolorbox}[breakable,colback=teal!4,colframe=teal!35!black,title={Before you move on}]
What does \hi{पहेली} mean? (\textquotedblleft{}A riddle\textquotedblright{}.) Say it once more.
\end{tcolorbox}
\section[\texorpdfstring{\hi{कहावत} (kahāvat)}{कहावत (kahāvat)}]{\hi{कहावत} (kahāvat) — a saying; a proverb}

\label{lesson:HI-C334-kahavat}



\begin{quote}
Before the new one: say the Hindi for a hermitage, then the Hindi for a riddle.
\end{quote}

\subsection*{You'll want to know: \hi{कहावत}}
\textbf{\hi{कहावत}} — \emph{kahāvat} — \textquotedblleft{}a saying; a proverb\textquotedblright{}.

\begin{quote}
\hi{दादी} \hi{हर} \hi{कहानी} \hi{के} \hi{बाद} \hi{एक} \hi{कहावत} \hi{कहती} \hi{थीं।}
\end{quote}

\begin{quote}
\emph{dādī har kahānī ke bād ek kahāvat kahtī thīṁ.}
\end{quote}

\begin{quote}
After every story my grandmother would say a proverb.
\end{quote}

Feminine.

\begin{cousinweb}
\textbf{\hi{कहावत}} (\emph{kahāvat}) is inherited from \emph{*kathāvārttā}, a form CDIAL reconstructs (CDIAL 2706), joining words for telling and for talk: a saying everyone repeats.
\end{cousinweb}

\begin{grammarlens}[title={What you've built}]
One more word for telling what happened.
\end{grammarlens}

\subsection*{Your turn}
\begin{itemize}
\raggedright
  \item \emph{Say it:} \emph{kahāvat}
  \item \emph{Say it:} \emph{kahāvat}, once more
  \item \emph{Say it:} \emph{dādī har kahānī ke bād ek kahāvat kahtī thīṁ.} — After every story my grandmother would say a proverb.
  \item \emph{Recall:} say the Hindi for a hermitage, then the Hindi for a riddle, then say \emph{kahāvat} again
\end{itemize}

\begin{tcolorbox}[breakable,colback=teal!4,colframe=teal!35!black,title={Before you move on}]
What does \hi{कहावत} mean? (\textquotedblleft{}A saying; a proverb\textquotedblright{}.) Say it once more.
\end{tcolorbox}
